\documentclass[10pt,a4paper]{amsart}
\usepackage[margin=19mm]{geometry}
\usepackage{amsmath,amssymb,mathtools}
\usepackage[scr=rsfs,cal=euler]{mathalfa}
\usepackage{xfrac}
\usepackage[unicode,hidelinks,bookmarks=false]{hyperref}
\newcommand{\ie}{i.e.\ }
\newcommand{\cf}{cf.\ }
\newcommand{\resp}{respectively\ }
\newcommand{\Subsection}{\subsection}
\providecommand{\nobreakdash}{\nobreak\hbox{-}}
\setlength{\emergencystretch}{2em}
\multlinegap0pt
\theoremstyle{plain}
\newtheorem{theorem}{Theorem}
\newtheorem{definition}[theorem]{Definition}
\newtheorem{prop}[theorem]{Proposition}
\newtheorem{lemma}[theorem]{Lemma}
\newtheorem{cor}[theorem]{Corollary}
\newtheorem{rem}[theorem]{Remark}
\newtheorem{ex}[theorem]{Example}
\renewcommand{\proof}{{\noindent \bf Proof. \ }}
\renewcommand{\proofname}{\noindent\bf Proof}
\newcommand{\bproof}{\noindent {\bf Proof:\ }}
\newcommand{\eproof}{\hfill \mbox{${\square}$}}
\newcommand{\R}{\mathbb{R}}
\newcommand{\dist}{\textrm{dist}}
\title[A Note on the formulation of the Neumann boundary condition for a nonlocal problem]{A Note on the formulation of the Neumann boundary condition for a nonlocal problem: the measure-preserving outer-collar involution (printed as Theorem 2.5)}
\author{Antonio Luiz Pereira}
\date{Source: arXiv:2509.16041v2 (CC BY 4.0)}
\begin{document}
\maketitle

\noindent\emph{Source and licence.} A Note on the formulation of the Neumann boundary condition for a nonlocal problem. Version arXiv:2509.16041v2 (CC BY 4.0), distributed under the Creative Commons Attribution 4.0 International licence, http://creativecommons.org/licenses/by/4.0/, as recorded in the arXiv licence metadata for this exact version and shown on the saved abstract page https://arxiv.org/abs/2509.16041v2. Full licence notice: licenses/arxiv-2509.16041-CC-BY-4.0.txt.\par\smallskip

\noindent\textit{Editorial scope.} Counted unit: the existence of the measure-preserving outer-collar involution for a bounded domain, printed in the article as Theorem 2.5 (source0.tex lines 390--400, one \texttt{theorem} environment with source label \texttt{area\_preserving}), delivered together with the article's complete original proof of it (source0.tex lines 402--453, the article's own run-in proof block). The article marks its proofs with the run-in macro \texttt{\textbackslash proof} and closes them with \texttt{\textbackslash eproof}, which prints a filled square; this package restates those two macros exactly as the article defines them and delivers the block from the run-in marker to the closing square with nothing added, so the proof reads exactly as it does in the article. No mathematics is written, repaired or re-derived, and no line of the statement or of the proof is omitted, substituted or re-flowed.

Required context retained uncounted: the article's normal-coordinates theorem with its label (source0.tex lines 343--374), which the counted statement and its proof both invoke, and which supplies the map $\Psi$, the nearest-point projection, the signed distance and the outward normal used throughout the counted proof. The two labelled displays the proof refers to, the normal-coordinate integration formula and the differential equation for the height function, are labelled inside the counted proof itself.

References. No citation command occurs in any retained line, so this package carries no bibliography block and no bibliography entry.

Editorial formatting only, in addition: an AMS \texttt{amsart} document, the same class the article itself uses, that restates the macro definitions the retained lines use (the article's own $\R$, $\dist$, its proof and end-of-proof markers) and the article's own theorem-environment declarations, keeping the article's convention that the lemma, corollary, remark and example environments share the theorem counter. Consequently the retained environments print here as Theorem 1 and Theorem 2 whereas the article prints them as Theorem 2.4 and Theorem 2.5, the article's own per-section numbering is not reproduced, and the equations of the retained spans are numbered anew in order of appearance; every cross-reference inside the retained text resolves to this wrapper's numbering, and that is the only change to the numbering. No figure, no image-inclusion line and no drawing code occur in any retained line, so no artwork is delivered or needed, although the article contains a figure elsewhere.

The article's own title line is longer than the arXiv record for this version: the source titles the note A note on the formulation of the Neumann boundary condition for a nonlocal diffusion problem with continuous kernel, while the record and the licence metadata give the title used here. The article has no separate class or style file; the package ships the article's concatenated source verbatim as \texttt{sources/185344/source0.tex}, and every delivered excerpt is an exact line range of that file. Not retained: the article's introduction, its other results, its examples and remarks, its final section and its bibliography; none of that material is used as a step of the delivered proof, which is complete as the author wrote it. No mathematical statement, hypothesis, estimate or conclusion has been rewritten, repaired or added, and printed slips, if any, are preserved literally.
\begin{theorem} \label{normal_neighborhood}  Let $ \Omega \subset \R^n$ be a {bounded} domain with a $\mathcal{C}^m$ boundary, $ m \geq 2$. There exists $r>0$ so that, if
  \begin{itemize}
  \item $B_r(\partial \Omega)= \{ x \in \R^n : \dist(x, \partial \Omega) < r  \}$,
 \item  $\pi(x) = \textrm{ the point of  }  \partial \Omega
 \textrm{ nearest to  } x  $,
 \item  $  t(x) =  \pm  \dist(x, \partial \Omega)$
  ($+$ outside, $-$ inside),
  \end{itemize}
then:
\begin{itemize}
\item  $ t(\cdot) : B_r(\partial \Omega) \to (-r,r) $ and
 $\pi(\cdot) :  B_r(\partial \Omega) \to \partial \Omega $ are
 well-defined,
 $\pi$ is  a $\mathcal{C}^{m-1} $ retraction onto
  $\partial \Omega$
and $t$ is of class $\mathcal{C}^{m} $;
 \item
$ x \mapsto  (\pi (x), t(x) ) :  B_r(\partial \Omega) \to \partial \Omega \times (-r,r)$  is a  $\mathcal{C}^{m-1} $
 diffeomorphism with inverse
$ \Psi(\xi, t) = \xi + t N(\xi)$,
 where $N(\xi) $ is the unique outward unit normal to $\partial \Omega$
  at $\xi$.
\end{itemize}
 \noindent Furthermore,
 $t(\cdot)$ is the unique solution of
 $| \nabla t(x)| = 1$ in   $B_r(\partial \Omega)$ with $t=0$ on
  $\partial \Omega$ and  $\frac{\partial t }{\partial  N } > 0 $ on  $\partial \Omega$.
  Extending the normal field $N$ to $B_r(\partial \Omega)$ by
$ N(\xi + t N(\xi) )= N(\xi)$, $-r <t < r $, we have
 $ N (x) = \nabla t(x)$  on $B_r (\partial \Omega).$ Also
 $ {\mathcal{S}} (x) = D N (x) = D^2 t(x)$, restricted to the tangent space at $ x  \in \partial \Omega$,  is the curvature {(shape operator)} of $\partial \Omega$.
\end{theorem}

\begin{theorem} \label{area_preserving} Let $ \Omega \subset \R^n$ be a {bounded} domain with a $\mathcal{C}^m$ boundary, $ m \geq 2$,
{and let $r>0$ be as in Theorem \ref{normal_neighborhood}. Then there exist $0 < \varepsilon < r$, $\varepsilon' > 0$ and a map
$\Phi : U_\varepsilon \to \R^n$, where $U_\varepsilon := \{x \in \Omega : \dist(x,\partial\Omega) < \varepsilon\}$,
satisfying the following properties:}
 \begin{enumerate}
  \item {$\Phi$ is a homeomorphism of $U_\varepsilon$ onto an open set $V \subset \Omega^c$}{, where $\Omega^c$ is the complement set of $\Omega$ in $\R^n$}, {and $V \supset \{ z \in \R^n \setminus \overline\Omega : \dist(z,\partial\Omega) < \varepsilon'\}$;}
  \item {$\Phi$ preserves Lebesgue measure: $\int_{U_\varepsilon} f(\Phi(x))\,dx = \int_V f(z)\,dz$ for every measurable $f \geq 0$;}
  \item {$\Phi$ is the restriction of an involution: the formula below defines $\Phi$ also on $V$, and $\Phi\circ\Phi = \mathrm{id}$ on $U_\varepsilon \cup V$;}
  \item {if $m \geq 3$, $\Phi$ is a $\mathcal{C}^{m-2}$ diffeomorphism of $U_\varepsilon$ onto $V$ and $\det D\Phi (x) = -1$ for all $x \in U_\varepsilon$.}
  \end{enumerate}
\end{theorem}

\noindent \proof
{Let $\Psi(y,s) = y + sN(y)$, $(y,s) \in \partial\Omega\times(-r,r)$, be the inverse of the normal coordinates given by Theorem \ref{normal_neighborhood}.
Let $\lambda_1(y), \dots, \lambda_{n-1}(y)$ be the eigenvalues of the curvature $\mathcal{S}(y) = DN(y)$ restricted to $T_y\partial\Omega$, and let
$v_1, \dots, v_{n-1}$ be a corresponding orthonormal basis of eigenvectors of $T_y\partial\Omega$.
Since $D\Psi_{(y,s)}(v_i,0) = v_i + s\,\mathcal{S}(y) v_i = (1+s\lambda_i)\, v_i$ and $D\Psi_{(y,s)}(0,1) = N(y)$, we have
\[
 \det D\Psi_{(y,s)} = \det\big[(1+s\lambda_1)v_1, \dots, (1+s\lambda_{n-1})v_{n-1}, N(y)\big] = p_y(s),
\]
where
\[ p_y(s) := \prod_{i=1}^{n-1}(1+s\lambda_i(y)) = \det\big(I + s\,\mathcal{S}(y)\big|_{T_y\partial\Omega}\big),\]
the determinant being computed with respect to the product of the surface measure $dS$ on $\partial\Omega$ and $ds$. Since $\Psi$ is a diffeomorphism, $p_y(s) \neq 0$ for $|s| < r$, and $p_y(0) = 1$, so $p_y(s) > 0$. Thus, for every measurable $f \geq 0$ supported in $B_r(\partial\Omega)$,
\begin{equation}\label{normal_coord_integral}
 \int_{\R^n} f(x)\, dx = \int_{\partial\Omega}\int_{-r}^{r} f\big(y + sN(y)\big)\, p_y(s)\, ds\, dS(y).
\end{equation}
The coefficients of the polynomial $p_y$ are the elementary symmetric functions of the $\lambda_i(y)$, which are continuous (indeed $\mathcal{C}^{m-2}$) in $y$.}

{For each $y \in \partial\Omega$ consider the o.d.e.}
\begin{equation}
 \label{diff_equation} \begin{cases}
 \dfrac{d\varphi}{ds} = -\dfrac{(1+ s \lambda_1(y))(1+ s \lambda_2(y))
 \cdots(1+ s \lambda_{n-1}(y))}{(1+ \varphi \lambda_1(y))(1+ \varphi \lambda_2(y))
 \cdots(1+ \varphi \lambda_{n-1}(y))} {= -\dfrac{p_y(s)}{p_y(\varphi)}},\\[2mm]
 \varphi(y,0) = 0.
\end{cases} \end{equation}
{This equation is separable: writing $P_y(s) := \int_0^s p_y(\sigma)\, d\sigma$, which is a strictly increasing function on $(-r,r)$, its solution is
\[ \varphi(y,s) = P_y^{-1}\big(-P_y(s)\big). \]
By continuity and compactness of $\partial\Omega$, there are $c, C > 0$ with $P_y(r/2) \geq c$ and $0 < p_y(s) \leq C$ for all $y \in \partial\Omega$, $|s| \leq r/2$; choosing $0 < \varepsilon < \min\{r/2, c/C\}$, we have $|P_y(s)| < c$ for $|s| < \varepsilon$, so that $\varphi(y,s)$ is well defined, continuous in $(y,s)$, $\mathcal{C}^1$ in $s$, and $\varphi(y,s) \in (-r/2,r/2)$ for $|s| < \varepsilon$.
Moreover, $\partial_s\varphi < 0$, so $\varphi(y,\cdot)$ maps $(-\varepsilon, 0)$ decreasingly onto $(0, \varphi(y,-\varepsilon))$, and $P_y(\varphi(y,\varphi(y,s))) = -P_y(\varphi(y,s)) = P_y(s)$, that is, $\varphi(y,\cdot)$ is an involution.
Note that $P_y(s)$ is the volume, per unit area of $\partial\Omega$ at $y$, of the layer between $\partial \Omega$ and $\partial\Omega + sN$; thus $\varphi(y,s)$ is the height of the outer layer with the same volume as the inner layer of depth $|s|$.}

{Define $\xi(y,s) := (y, \varphi(y,s))$ and
  \begin{equation*}
   \Phi(x) := \Psi\circ\xi\circ\Psi^{-1}(x) = \pi(x) + \varphi( \pi(x), t(x))\, N(\pi(x)), \qquad x \in U_\varepsilon,
  \end{equation*}
and let $V := \Psi\big(\{(y,\sigma) : y \in\partial\Omega,\ 0 < \sigma < \varphi(y,-\varepsilon)\}\big)$.
Since $\xi$ is a homeomorphism of $\partial\Omega\times(-\varepsilon,0)$ onto $\{(y,\sigma): 0<\sigma<\varphi(y,-\varepsilon)\}$ (with inverse given by $\xi$ itself), $\Phi$ is a homeomorphism of $U_\varepsilon$ onto $V$, and $\Phi\circ\Phi = \mathrm{id}$ on $U_\varepsilon\cup V$. Since $t > 0$ on $V$, $V \subset \Omega^c$, and $V \supset \{z : 0 < t(z) < \varepsilon'\}$ with $\varepsilon' := \min_{y\in\partial\Omega} \varphi(y,-\varepsilon) > 0$. This proves 1 and 3.}

{To prove 2, let $f \geq 0$ be measurable. By \eqref{normal_coord_integral} and the change of variables $\sigma = \varphi(y,s)$ for fixed $y$, for which $p_y(s)\,ds = -p_y(\sigma)\,d\sigma$ by \eqref{diff_equation},
\begin{align*}
\int_{U_\varepsilon} f(\Phi(x))\, dx
 & = \int_{\partial\Omega}\int_{-\varepsilon}^{0} f\big(y + \varphi(y,s)N(y)\big)\, p_y(s)\, ds\, dS(y) \\
 & = \int_{\partial\Omega}\int_{0}^{\varphi(y,-\varepsilon)} f\big(y + \sigma N(y)\big)\, p_y(\sigma)\, d\sigma\, dS(y)
 = \int_V f(z)\, dz .
\end{align*}
Finally, if $m \geq 3$, then $\varphi$ is $\mathcal{C}^{m-2}$ in $(y,s)$ and so is $\Phi$. The matrix of $D\xi(y,s)$ is block triangular,
$D\xi(y,s)(v,0) = (v, D_y\varphi(y,s)\,v)$, $D\xi(y,s)(0,1) = (0, \partial_s\varphi(y,s))$, so $\det D\xi = \partial_s\varphi$, and
\[
 \det D\Phi(x) = \det D\Psi_{\xi(y,s)}\cdot \partial_s\varphi(y,s)\cdot \big(\det D\Psi_{(y,s)}\big)^{-1}
 = p_y(\varphi)\cdot\Big(-\frac{p_y(s)}{p_y(\varphi)}\Big)\cdot \frac{1}{p_y(s)} = -1,
\]
where $(y,s) = \Psi^{-1}(x)$ and $\varphi = \varphi(y,s)$. This proves 4.}
  \eproof

\end{document}
